\section{Polynomial coefficients as ordered gap sums}

\begin{theorem}[Ordered gap expansion of a coefficient]
    \label{thm:polynomial_ordered_gaps}
    \lean{Polynomial.coeff_pow_C_add_C_mul_X_eq_sum_gaps}
    \leanok
    Let $\mathcal R$ be a semiring, let $D,T\in\mathcal R$, and let
    $m,k\ge0$ be integers. Write $[X^k]P$ for the coefficient of degree $k$
    in a polynomial $P\in\mathcal R[X]$, where $X$ is central. Then
    \begin{align*}
        [X^k](D+TX)^{m+k}
        =\sum_{\substack{n_0,\ldots,n_k\ge0\\n_0+\cdots+n_k=m}}
            D^{n_0}TD^{n_1}\cdots TD^{n_k}.
    \end{align*}
    The factors retain their displayed order. At $k=0$ the sum is $D^m$;
    at $m=0$ it is $T^k$.
    This identity gives the ordered collection step in the proof of
    \cite[Lemma~4.2]{Arad2013AreaLaw}, before expanding the two commuting
    contributions within each gap.
\end{theorem}
\begin{proof}
    \leanok
    Let $F_{m,k}$ be the coefficient on the left and $G_{m,k}$ the gap sum
    on the right. Multiplying the polynomial power on the left by $D+TX$
    gives
    \begin{align*}
        F_{m+1,k+1}=D F_{m,k+1}+T F_{m+1,k}.
    \end{align*}
    For the gap sum, removing one $D$ from a positive first gap gives the
    first term, while removing the first $T$ when its preceding gap is zero
    gives the second. Thus
    \begin{align*}
        G_{m+1,k+1}=D G_{m,k+1}+T G_{m+1,k}.
    \end{align*}
    The boundary identities are
    \begin{align*}
        F_{m,0}=G_{m,0}=D^m, \\
        F_{0,k}=G_{0,k}=T^k.
    \end{align*}
    Induction on $m$, followed by induction on $k$, proves $F_{m,k}=G_{m,k}$.
    The list of gap tuples has no repetitions, so its sum equals the sum
    over the corresponding finite set of tuples.
\end{proof}
